\appendix
\chapter*{Appendix}
\addcontentsline{toc}{chapter}{Appendix}
\label{app:supplementary-materials}

\setcounter{section}{0}
\renewcommand{\thesection}{A.\arabic{section}}
\renewcommand{\thesubsection}{A.\arabic{section}.\arabic{subsection}}
\renewcommand{\thesubsubsection}{A.\arabic{section}.\arabic{subsection}.\arabic{subsubsection}}

This appendix documents the AI-assisted material used during thesis preparation
and provides supplementary evaluation and traceability information. In line
with the applicable thesis guidelines, large software artefacts and
machine-readable logs are retained in the electronic submission rather than
reproduced in full in the printed appendix.

\section{AI Usage Documentation}
\label{app:ai-usage}

Generative AI assistance was used in a limited and supportive capacity during
document preparation. The support was restricted to language refinement,
grammar correction, structure checking, formatting support, LaTeX assistance,
and consistency review. It was not used to generate experimental data,
measurement values, evaluation outcomes, system requirements, implementation
logic, or scientific conclusions.

All AI-assisted text and formatting suggestions were reviewed by the author
before use. The author retained responsibility for the technical content,
interpretation of results, and final thesis wording.

\section{Prompt Log and Supporting Files}
\label{app:prompt-log}

A machine-readable prompt log is retained in the electronic submission as
\texttt{appendix/ai\_prompt\_log.txt}. The log is written in a
YAML-compatible plain-text structure and records the retained AI usage prompts,
the AI tool, provider, model family, purpose, and associated files where
applicable. Supporting thesis source files referenced by those entries are
included in the submission package.

\section{Supplementary Evaluation Documentation}
\label{app:supplementary-evaluation}

Chapter~\ref{ch:evaluation-validation} presents the simulation-based
evaluation of the implemented system. The evaluation is based on repeated
manual-versus-automated workflow trials and automated scenario trials covering
nominal Side~A bending, opposite-side Side~E bending, geometry variation,
workpiece position variation, regrip-required workflows, and failure-prone
conditions. Table~\ref{tab:manual-vs-auto-trial-structure} summarises the
evaluation groups and the corresponding measured quantities.

\begin{table}[htbp]
\centering
\normalsize
\setlength{\tabcolsep}{3pt}
\renewcommand{\arraystretch}{1.12}
\begin{tabularx}{\textwidth}{
L{3.0cm}
L{3.4cm}
Y
L{3.0cm}
}
\hline
\textbf{Evaluation Group} &
\textbf{Scenario Coverage} &
\textbf{Purpose} &
\textbf{Measured Metrics} \\
\hline

Manual workflow trials &
Side~A, Side~E, geometry variation, and workpiece position variation &
Quantify manual preparation effort, simulated execution time, and total process time for comparison with automated generation. &
Preparation, execution, total time, regrip count \\
\hline

Automated workflow trials &
Side~A, Side~E, geometry variation, and workpiece position variation &
Quantify automated preparation effort, simulated execution time, and total process time under the same workflow classes. &
Preparation, execution, total time, regrip count \\
\hline

Automated feasibility trials &
Nominal, orientation variation, workspace-constrained, regrip-required, opposite-side, and failure-prone scenarios &
Verify scenario-wise completion, controlled-stop behaviour, planning time, and robustness limits of the generated workflow. &
Success rate, failure rate, planning time, controlled stop \\

\hline
\end{tabularx}
\caption{Evaluation trial structure used for the final simulation-based validation.}
\label{tab:manual-vs-auto-trial-structure}
\end{table}

\FloatBarrier


The table is included to make the measurement structure transparent without
duplicating the full machine-readable trial records in the printed appendix.
The complete numerical analysis and interpretation are provided in
Chapter~\ref{ch:evaluation-validation}.

\section{Requirement Traceability and Compliance Tables}
\label{app:req-traceability-compliance}

The following tables provide the detailed requirement traceability and
implementation compliance mappings referenced from
Chapters~\ref{ch:system-concept-methodology} to~\ref{ch:implementation}.
They are placed in the appendix to keep the main thesis flow focused while
preserving the audit trail between requirements, implementation modules, and
validation activities. The mappings are aligned with the scenario-based
verification presented in Chapter~\ref{ch:evaluation-validation}. Tables~\ref{tab:req-traceability-matrix}
and~\ref{tab:req-traceability-matrix-nfr-sr} provide the traceability mapping,
while Tables~\ref{tab:req-compliance} and~\ref{tab:req-compliance-nfr-sr}
summarise the corresponding implementation compliance mapping.

\begin{landscape}

\begin{table}[p]
\centering
\normalsize
\setlength{\tabcolsep}{3pt}
\renewcommand{\arraystretch}{1.08}

\begin{tabular}{
L{1.8cm}
L{6.0cm}
L{9.0cm}
L{5.5cm}
}
\hline
\textbf{Req.\ ID} &
\textbf{Requirement} &
\textbf{Implementation} &
\textbf{Verification} \\
\hline

FR-00 &
Generate execution-ready drag\&bot programs from geometry, configuration data, and bend-side input &
Integrated planning pipeline for configuration processing, orientation selection, grip planning, trajectory generation, and execution logic &
Manual-versus-automated workflow comparison and automated scenario validation in Chapter~\ref{ch:evaluation-validation} \\
\hline

FR-01 &
Detect unreachable or infeasible target poses before command dispatch &
Reachability checks and admissibility validation &
Workspace-constrained and failure-prone scenario validation with controlled-stop classification \\
\hline

FR-02 &
Bend-side selection before automated program generation or execution &
Operator dialog and configuration input logic &
Side~A direct-bending and Side~E opposite-side workflow validation \\
\hline

FR-03 &
Automatically insert a regrip operation when the next target pose is not directly reachable &
Regrip cycle logic with \acs{io} and attach-detach control &
Regrip-required and Side~E opposite-side scenario validation \\
\hline

FR-04 &
Automatic orientation validation and correction for changed initial workpiece orientation &
Orientation conversion, rotation normalisation, and strategy selection &
Orientation-variation scenario validation \\
\hline

FR-05 &
Automatically generate flip and reorientation workflow for opposite-side bending &
Opposite-side workflow with sheet flip, regrip, revalidation, and controlled-stop handling if recovery fails &
Side~E opposite-side workflow validation \\
\hline

FR-06 &
Determine grip positions that satisfy edge-clearance, hole-avoidance, and holding-force checks &
Grip planning with edge-clearance, spacing, and vacuum-force checks &
\acs{tc}-06 grip-planning test and functional validation of valid and rejected grips \\
\hline

FR-07 &
Detect unsafe conditions and replan or stop execution before unsafe continuation &
Safety validation, pre-execution checks, replanning, and controlled stop logic &
Failure-prone scenario validation and controlled-stop classification \\

\hline
\end{tabular}

\caption[Requirement Traceability Matrix -- functional requirements]{Requirement Traceability Matrix -- functional requirements (FR-00 to FR-07) mapped to implementation modules and verification scenarios.}
\label{tab:req-traceability-matrix}

\end{table}


\begin{table}[p]
\centering
\normalsize
\setlength{\tabcolsep}{3pt}
\renewcommand{\arraystretch}{1.08}

\begin{tabular}{
L{1.8cm}
L{6.0cm}
L{9.0cm}
L{5.5cm}
}
\hline
\textbf{Req.\ ID} &
\textbf{Requirement} &
\textbf{Implementation} &
\textbf{Verification} \\
\hline

NFR-01 &
Limit operator interaction to configuration or input selection and execution confirmation &
Bend-side selection and execution confirmation only &
Usability validation and workflow comparison without manual waypoint or block programming \\
\hline

NFR-02 &
Configure workpiece geometry, process settings, and safety parameters without source-code modification &
Configuration reader and \acs{yaml}-based process configuration &
Geometry-variation and position-change scenario validation through configuration or operator input updates \\
\hline

NFR-03 &
Provide warning or error feedback before continuation for unsafe or infeasible conditions &
Warning generation and operator notification logic &
Warning message validation \\
\hline

NFR-04 &
Deterministic execution sequence for a given configuration &
Deterministic execution loop with automatic waypoint generation and logging &
Repeated automated scenario validation using completion, failure, or controlled-stop classifications \\
\hline

NFR-05 &
Check fixture pose deviations against defined positional and angular tolerance limits &
Configurable pose-deviation threshold and perturbation handling &
Fixture perturbation and constrained-placement validation using 5~mm and 0{,}10~rad thresholds \\
\hline

\acs{sr}-01 &
Connect to and validate generated programs in the configured drag\&bot virtual execution environment &
Runtime client connection to the configured drag\&bot environment &
Platform integration check \\
\hline

\acs{sr}-02 &
Route planning-to-runtime communication through the AnyRobot WebSocket interface &
WebSocket-based communication layer for structured command exchange &
Interface integration check \\
\hline

\acs{sr}-03 &
Generate all motion targets using the configured \ac{tcp} identifier \texttt{project\_3\_tcp} &
\acs{tcp}-based transformations in trajectory generation and execution &
\acs{tcp} configuration check \\

\hline
\end{tabular}

\caption[Requirement Traceability Matrix -- non-functional and system requirements]{Requirement Traceability Matrix -- non-functional requirements (NFR-01 to NFR-05) and system requirements (SR-01 to SR-03) mapped to implementation modules and verification scenarios.}
\label{tab:req-traceability-matrix-nfr-sr}

\end{table}

\end{landscape}



\begin{table}[htbp]
\centering
\caption{Requirement compliance matrix linking system requirements to implementation components}
\label{tab:req_compliance}
\small
\begin{tabularx}{\textwidth}{
L{2.5cm}
L{5.5cm}
L{6.5cm}
}
\toprule
\textbf{Requirement ID} & 
\textbf{Requirement Description} & 
\textbf{Implementation Mapping} \\
\midrule
 
FR-00 & Overall system shall automate press brake bending process workflow & 
Achieved through the integrated planning pipeline, including orientation strategy (Section~\ref{sec:workpiece_orientation}), grip planning (Section~\ref{sec:grip_planning}), regrip planning (Section~\ref{sec:regrip}), and trajectory generation (Section~\ref{sec:trajectory}) \\
 
FR-01 & Automatic generation of robot programs from geometric input & 
Implemented through the structured planning pipeline combining orientation determination, grip planning, and trajectory generation (Sections~\ref{sec:workpiece_orientation}, \ref{sec:grip_planning}, \ref{sec:trajectory}) \\
 
FR-02 & Support for multiple workpiece orientations and configurations & 
Handled through the three-level orientation strategy with admissibility validation and fallback mechanism (Section~\ref{sec:workpiece_orientation}) \\
 
FR-03 & Automatic grip position determination & 
Achieved through constraint-based grip planning with geometric filtering and force validation (Section~\ref{sec:grip_planning}) \\
 
FR-04 & Handling of regrip scenarios for unreachable configurations & 
Implemented via dynamic regrip planning and execution strategy based on kinematic feasibility evaluation (Section~\ref{sec:regrip}) \\
 
FR-05 & Support for opposite-side bending operations & 
Enabled through the flip workflow and reuse of the planning pipeline for secondary operations (Section~\ref{sec:opposite_side}) \\
 
FR-06 & Seamless communication with robot runtime & 
Achieved using the AnyRobot Interface with authenticated WebSocket communication (Section~\ref{sec:communication_stack}) \\
 
FR-07 & Pre-execution safety validation before motion execution & 
Implemented through fixture pose validation, operator confirmation, and runtime monitoring mechanisms (Section~\ref{Pre-Execution_Safety_Validation}) \\
 
NFR-01 & Deterministic and repeatable system behaviour & 
Ensured through a structured modular pipeline and fixed execution order across system layers (Sections~5.2–5.10) \\
 
NFR-02 & Configurability without source code modification & 
Achieved through centralized YAML configuration architecture and parameter-driven system design (Section~\ref{Configuration Architecture }) \\
 
NFR-03 & Safe system operation with operator control & 
Enforced through blocking operator confirmation and structured safety validation (Section~\ref{Pre-Execution_Safety_Validation}) \\
 
NFR-04 & Early detection of invalid configurations & 
Ensured through configuration validation and fixture deviation checks prior to execution (Section~\ref{Pre-Execution_Safety_Validation}) \\
 
NFR-05 & Robust handling of runtime errors & 
Implemented through continuous runtime monitoring using the command result watcher and HI notifications (Section~\ref{Pre-Execution_Safety_Validation}) \\
 
\bottomrule
\end{tabularx}
\end{table}



\section{Hardware Implementation Report}
\label{app:hardware-implementation-report}

The implemented system was evaluated in a configured virtual press-brake cell
rather than on a physical production installation. This section records the
hardware and platform configuration used for the simulation-based
implementation so that the scope of the results remains transparent and
consistent with Chapter~\ref{ch:evaluation-validation}.

\begin{table}[htbp]
\centering
\normalsize
\setlength{\tabcolsep}{4pt}
\renewcommand{\arraystretch}{1.12}
\begin{tabularx}{\textwidth}{L{4.0cm} Y}
\hline
\textbf{Configuration item} & \textbf{Implementation report} \\
\hline
drag\&bot platform &
drag\&bot runtime package 3.6.0, used as the virtual robot-programming and
execution environment for the thesis implementation. \\
\hline
Runtime interface &
AnyRobot WebSocket-based interface used for command exchange between the
planning software and the drag\&bot runtime. \\
\hline
Simulation setup &
Virtual press-brake cell containing the UR5e robot model, press brake,
orientation table, flip station, vacuum gripper, and tracked sheet object. \\
\hline
Robot configuration &
UR5e manipulator model with six degrees of freedom, used for
all simulated handling and bending-related motion sequences. \\
\hline
\ac{tcp} definition &
Configured TCP \texttt{project\_3\_tcp} with offsets
$x = 0{,}000$~m, $y = 0{,}034$~m, $z = 0{,}194$~m,
$RZ = -90{,}0^\circ$, $RY = -75{,}0^\circ$, and $RX = 0{,}0^\circ$. \\
\hline
Gripper setup &
Two-cup vacuum gripper represented in the simulation by digital input-output
control, vacuum activation and release states, and kinematic object
attachment-detachment logic. \\
\hline
System configuration &
Workpiece dimensions, bend-side selection, safety limits, admissible workspace
zones, gripper parameters, and runtime settings are loaded from configuration
data before execution. \\
\hline
\end{tabularx}
\caption{Hardware and platform configuration used for the simulation-based implementation.}
\label{tab:hardware-implementation-report}
\end{table}


This configuration supports the reported simulation-based validation, but it
does not replace physical commissioning. Physical robot validation, bend-angle
measurement, gripper slip testing, forming-force assessment, and certified
safety verification remain outside the scope of the present thesis.
